\scanpage{380}
\mspage{296}
% archive sheet number: 379

cofinale et g\textsuperscript{r}de devant $\pi$, \uline{pourvu que}
$\operatorname{card} A\leq\pi$, et que l'on \deleted{suppose}\inserted{postule} une
propriété de stabilité de $I$ et des
$J_\alpha$ dans $I$, pour les limites inductives
de cardinal $\leq\pi$ et g\textsuperscript{r}des devant un
cardinal convenable $\pi'_0<\pi$ donné.

4.13.10.1. Notons d'ailleurs que si
\deleted{\uline{Démonstration \unsure{du lemme|des lemmes}}}
$f:J\unsure{\text{---}|\longrightarrow}I$ est un foncteur \inserted{cofinal} \uline{plein\unsure{mt}} \uline{fidèle},
avec $I$ g\textsuperscript{r}de devant $\pi$, alors $J$ est
automatiquement g\textsuperscript{r}de devant $\pi$. Cela
résulte p. ex. de 4.9.6.3 $2^\circ$ (p. \unsure{225|235})
par le critère $C_\pi$ (vu que pour $f$ pl fid,
la condition $\unsure{\mathrm E.3_{1\pi}|\mathrm E.3_{17}}$ est tautologiquement
satisfaite). On peut aussi le voir ainsi.
Soit $B$ un ens. avec $\operatorname{card}B\leq\pi$,
il faut prouver que $\deleted{\unsure{B}}\unsure{J|B|I}\unsure{\text{---}|\longrightarrow}J^B$ est cofinal,
or on a

\diagram{\begin{tikzcd}[ampersand replacement=\&]
J \arrow[r,dash,"{\mathrm{cof.}?}"] \arrow[d,"{\mathrm{cof.}}"'] \& J^B \arrow[d,"{\mathrm{cof.}}"]\\
I \arrow[r,dash,"{\mathrm{cof.}}"] \& I^B \mathpunct{,}
\end{tikzcd}}

où trois des quatre flèches du carré
sont cofinales, $J^B\unsure{\text{---}|\longrightarrow}I^B$ l'étant comme
produit de foncteurs cofinaux entre
catégories \uline{filtrantes} (cf 4.9.6.5, page 227)
-- on sait en effet déjà que $J$ est filtrante.
